\frfilename{mt232.tex}
\versiondate{5.6.02}
\copyrightdate{1994}

\def\chaptername{Teorema Radon-Nikod\'ym}
\def\sectionname{Teorema Radon-Nikod\'ym}

\newsection{232}

Sekarang saya sampai pada teorema utama bab ini, salah satu hasil sentral
teori ukuran, yang menghubungkan fungsional aditif terhitung dengan integral
tak tentu. Tujuannya ialah memberikan deskripsi lengkap mengenai fungsional
yang dapat muncul sebagai integral tak tentu dari fungsi terintegralkan
(232E). Fungsional tersebut dapat dicirikan sebagai fungsional aditif yang
`kontinu sejati' (232Ab). Konsep yang lebih lazim digunakan, dan memadai
dalam banyak kasus, ialah fungsional aditif `kontinu mutlak' (232Aa); saya
menggunakan beberapa paragraf pertama (232B-232D) untuk membahas fakta-fakta
elementer mengenai fungsional kontinu sejati dan kontinu mutlak. Bagian ini
saya akhiri dengan pembahasan mengenai penguraian fungsional aditif terhitung
umum (232I).

\leader{232A}{Fungsional kontinu mutlak} Misalkan $(X,\Sigma,\mu)$ suatu
ruang ukur dan $\nu:\Sigma\to\Bbb R$ suatu fungsional aditif hingga.

\header{232Aa}{\bf (a)} $\nu$ bersifat {\bf kontinu mutlak} terhadap
$\mu$ (kadang-kadang ditulis `$\nu\ll\mu$') jika untuk setiap
$\epsilon>0$ terdapat $\delta>0$ sedemikian sehingga
$|\nu E|\le\epsilon$ setiap kali $E\in\Sigma$ dan $\mu E\le\delta$.


\header{232Ab}{\bf (b)} $\nu$ bersifat {\bf kontinu sejati} terhadap
$\mu$ jika untuk setiap $\epsilon>0$ terdapat $E\in\Sigma$ dan
$\delta>0$ sedemikian sehingga $\mu E$ berhingga dan
$|\nu F|\le\epsilon$ setiap kali $F\in\Sigma$ dan
$\mu(E\cap F)\le\delta$.


\header{232Ac}{\bf (c)}\cmmnt{ Untuk rujukan, saya tambahkan definisi
lain di sini.} Jika $\nu$ aditif terhitung, maka ia bersifat
{\bf singular} terhadap $\mu$ jika terdapat himpunan $F\in\Sigma$
sedemikian sehingga $\mu F=0$ dan $\nu E=0$ setiap kali $E\in\Sigma$
dan $E\subseteq X\setminus F$.


\leader{232B}{Proposisi} Misalkan $(X,\Sigma,\mu)$ suatu ruang ukur dan
$\nu:\Sigma\to\Bbb R$ suatu fungsional aditif hingga.

(a) Jika $\nu$ aditif terhitung, maka ia kontinu mutlak terhadap
$\mu$ jika dan hanya jika $\nu E=0$ setiap kali $\mu E=0$.

(b) $\nu$ kontinu sejati terhadap $\mu$ jika dan hanya jika
($\alpha$) ia aditif terhitung, ($\beta$) ia kontinu mutlak terhadap
$\mu$, dan ($\gamma$) setiap kali $E\in\Sigma$ dan $\nu E\ne 0$ terdapat
$F\in\Sigma$ sedemikian sehingga $\mu F<\infty$ dan
$\nu(E\cap F)\ne 0$.

(c) Jika $(X,\Sigma,\mu)$ $\sigma$-hingga, maka $\nu$ kontinu sejati
terhadap $\mu$ jika dan hanya jika ia aditif terhitung dan kontinu
mutlak terhadap $\mu$.

(d) Jika $(X,\Sigma,\mu)$ hingga total, maka $\nu$ kontinu sejati
terhadap $\mu$ jika dan hanya jika ia kontinu mutlak terhadap $\mu$.

\proof{{\bf (a)(i)} Jika $\nu$ kontinu mutlak terhadap $\mu$ dan
$\mu E=0$, maka $\mu E\le\delta$ untuk setiap $\delta>0$, sehingga
$|\nu E|\le\epsilon$ untuk setiap $\epsilon>0$ dan $\nu E=0$.

\medskip

\quad{\bf (ii)} \Quer\ Andaikan, jika mungkin, bahwa $\nu E=0$ setiap
kali $\mu E=0$, tetapi $\nu$ tidak kontinu mutlak. Maka terdapat
$\epsilon>0$ sedemikian sehingga untuk setiap $\delta>0$ terdapat
$E\in\Sigma$ dengan $\mu E\le\delta$ tetapi
$|\nu E|\ge\epsilon$. Untuk setiap $n\in\Bbb N$ kita dapat memilih
$F_n\in\Sigma$ sedemikian sehingga $\mu F_n\le 2^{-n}$ dan
$|\nu F_n|\ge\epsilon$. Tinjau
$F=\bigcap_{n\in\Bbb N}\bigcup_{k\ge n}F_k$. Maka

\Centerline{$\mu F\le\inf_{n\in\Bbb N}\mu(\bigcup_{k\ge n}F_k)
\le\inf_{n\in\Bbb N}\sum_{k=n}^{\infty}2^{-k}=0$,}

\noindent sehingga $\mu F=0$.

Sekarang ingat bahwa menurut 231Eb terdapat $H\in\Sigma$ sedemikian
sehingga $\nu G\ge 0$ jika $G\in\Sigma$ dan $G\subseteq H$, serta
$\nu G\le 0$ jika $G\in\Sigma$ dan $G\cap H=\emptyset$. Seperti dalam
231F, tetapkan $\nu_1 G=\nu(G\cap H)$ dan
$\nu_2G=-\nu(G\setminus H)$ untuk $G\in\Sigma$, sehingga $\nu_1$ dan
$\nu_2$ merupakan ukuran hingga total, dan $\nu_1F=\nu_2F=0$ karena
$\mu(F\cap H)=\mu(F\setminus H)=0$. Akibatnya

\Centerline{$0=\nu_iF=\lim_{n\to\infty}\nu_i(\bigcup_{m\ge n}F_m)
\ge\limsup_{n\to\infty}\nu_iF_n$}

\noindent untuk kedua nilai $i$, dan

\Centerline{$0=\lim_{n\to\infty}(\nu_1F_n+\nu_2F_n)
\ge\liminf_{n\to\infty}|\nu F_n|\ge\epsilon>0$,}

\noindent yang mustahil.\ \Bang

\medskip

{\bf (b)(i)} Andaikan $\nu$ kontinu sejati terhadap $\mu$. Dari
definisi-definisi tersebut jelas bahwa $\nu$ kontinu mutlak terhadap
$\mu$. Jika $\nu E\ne 0$, mesti terdapat $F$ yang berukuran hingga
sedemikian sehingga $|\nu G|<|\nu E|$ setiap kali
$G\cap F=\emptyset$, sehingga $|\nu(E\setminus F)|<|\nu E|$ dan
$\nu(E\cap F)\ne 0$. Ini menangani syarat ($\beta$) dan ($\gamma$).

Untuk memeriksa bahwa $\nu$ aditif terhitung, misalkan
$\sequencen{E_n}$ suatu barisan saling lepas dalam $\Sigma$ dengan
gabungan $E$, dan misalkan $\epsilon>0$. Ambil $\delta>0$ dan
$F\in\Sigma$ sedemikian sehingga $\mu F<\infty$ dan
$|\nu G|\le\epsilon$ setiap kali $G\in\Sigma$ dan
$\mu(F\cap G)\le\delta$. Maka

\Centerline{$\sum_{n=0}^{\infty}\mu(E_n\cap F)\le\mu F<\infty$,}

\noindent sehingga terdapat $n\in\Bbb N$ sedemikian sehingga
$\sum_{i=n}^{\infty}\mu(E_i\cap F)\le\delta$. Ambil sembarang $m\ge n$
dan tinjau $E^*_m=\bigcup_{i\le m}E_i$. Kita mempunyai

%$$\eqalign{|\nu E-\sum_{i=0}^m\nu E_i|
%&=|\nu E-\nu E^*_m|\cr
%&=|\nu(E\setminus E^*_m)|\cr
%&\le\epsilon\cr}$$

\Centerline{$|\nu E-\sum_{i=0}^m\nu E_i|
=|\nu E-\nu E^*_m|
=|\nu(E\setminus E^*_m)|
\le\epsilon$,}


\noindent karena

\Centerline{$\mu(F\cap E\setminus E^*_m)
=\sum_{i=m+1}^{\infty}\mu(F\cap E_i)
\le\delta$.}

\noindent Karena $\epsilon$ sembarang,

\Centerline{$\nu E=\sum_{i=0}^{\infty}\nu E_i$;}

\noindent karena $\sequencen{E_n}$ sembarang, $\nu$ aditif terhitung.

\medskip

\quad{\bf (ii)} Sekarang andaikan $\nu$ memenuhi ketiga syarat tersebut.
Menurut 231F, $\nu$ dapat dinyatakan sebagai selisih dua fungsional aditif
terhitung nonnegatif $\nu_1$, $\nu_2$; tetapkan
$\nuprime=\nu_1+\nu_2$, sehingga $\nuprime$ merupakan fungsional aditif
terhitung nonnegatif dan $|\nu F|\le\nuprime F$ untuk setiap
$F\in\Sigma$. Tetapkan

\Centerline{$\gamma=\sup\{\nuprime F:F\in\Sigma,\,\mu
F<\infty\}\le\nuprime
X<\infty$,}

\noindent dan pilih suatu barisan $\sequencen{F_n}$ yang terdiri atas
himpunan-himpunan berukuran hingga sedemikian sehingga
$\lim_{n\to\infty}\nuprime F_n=\gamma$; tetapkan
$F^*=\bigcup_{n\in\Bbb N}F_n$.
Jika $G\in\Sigma$ dan $G\cap F^*=\emptyset$, maka $\nu G=0$.
\Prf\Quer\
Jika tidak, menurut syarat ($\gamma$), terdapat $F\in\Sigma$
sedemikian sehingga $\mu F<\infty$ dan $\nu(G\cap F)\ne 0$. Maka

\Centerline{$\nuprime(F\setminus F^*)\ge\nuprime(F\cap G)\ge|\nu(F\cap
G)|>0$,}

\noindent dan mesti terdapat $n\in\Bbb N$ sedemikian sehingga

\Centerline{$\gamma<\nuprime F_n+\nuprime(F\setminus
F^*)=\nuprime(F_n\cup(F\setminus
F^*))\le\nuprime(F\cup F_n)\le\gamma$}

\noindent karena $\mu(F\cup F_n)<\infty$; tetapi ini mustahil.
\Bang\Qed

Dengan menetapkan $F_n^*=\bigcup_{k\le n}F_k$ untuk setiap $n$, kita
memperoleh $\lim_{n\to\infty}\nuprime(F^*\setminus F^*_n)=0$.
Ambil sembarang $\epsilon>0$, dan (dengan menggunakan syarat ($\beta$))
pilih $\delta>0$ sedemikian sehingga $|\nu E|\le{1\over 2}\epsilon$
setiap kali $\mu E\le\delta$. Pilih $n$ sedemikian sehingga
$\nuprime(F^*\setminus F_n^*)\le\bover12\epsilon$. Sekarang, jika
$F\in\Sigma$ dan $\mu(F\cap F^*_n)\le\delta$, maka

$$\eqalign{|\nu F|
&\le|\nu(F\cap F^*_n)|
   +|\nu(F\cap F^*\setminus F^*_n)|
   +|\nu(F\setminus F^*)|\cr
&\le\Bover12\epsilon+\nuprime(F\cap F^*\setminus F^*_n)+0\cr
&\le\Bover12\epsilon+\nuprime(F^*\setminus F^*_n)
\le\Bover12\epsilon+\Bover12\epsilon
=\epsilon.\cr}$$

\noindent Dan $\mu F^*_n<\infty$. Karena $\epsilon$ sembarang, $\nu$
kontinu sejati.

\medskip

{\bf (c)} Sekarang andaikan $(X,\Sigma,\mu)$ $\sigma$-hingga dan $\nu$
aditif terhitung serta kontinu mutlak terhadap $\mu$. Misalkan
$\sequencen{X_n}$ suatu barisan tak menurun yang terdiri atas himpunan
berukuran hingga dan meliputi $X$ (211D). Jika $\nu E\ne 0$, maka
$\lim_{n\to\infty}\nu(E\cap X_n)\ne 0$, sehingga
$\nu(E\cap X_n)\ne 0$ untuk suatu $n$. Ini menunjukkan bahwa $\nu$
memenuhi syarat ($\gamma$) dalam (b), sehingga kontinu sejati.

Tentu saja, implikasi sebaliknya sudah tercakup dalam (b).

\medskip

{\bf (d)} Terakhir, andaikan $\mu X<\infty$ dan $\nu$ kontinu mutlak
terhadap $\mu$. Maka ia mesti kontinu sejati, karena kita dapat
mengambil $F=X$ dalam definisi 232Ab.
}%end of proof of 232B

\leader{232C}{Lema} Misalkan $(X,\Sigma,\mu)$ suatu ruang ukur, dan
$\nu$, $\nuprime$ dua fungsional aditif terhitung pada $\Sigma$ yang
kontinu sejati terhadap $\mu$. Ambil $c\in\Bbb R$ dan $H\in\Sigma$,
serta tetapkan $\nu_HE=\nu(E\cap H)$\cmmnt{, seperti dalam 231De.}
Maka $\nu+\nuprime$, $c\nu$, dan $\nu_H$ semuanya kontinu sejati terhadap
$\mu$, dan $\nu$ dapat dinyatakan sebagai selisih fungsional-fungsional
aditif terhitung nonnegatif yang kontinu sejati terhadap $\mu$.

\proof{ Misalkan $\epsilon>0$. Tetapkan $\eta=\epsilon/(2+|c|)>0$.
Maka terdapat $\delta$, $\delta'>0$ serta $E$, $E'\in\Sigma$ sedemikian
sehingga $\mu E<\infty$, $\mu E'<\infty$, dan $|\nu F|\le\eta$ setiap
kali $\mu(F\cap E)\le\delta$, serta $|\nuprime F|\le\eta$ setiap kali
$\mu(F\cap E')\le\delta'$. Tetapkan
$\delta^*=\min(\delta,\delta')>0$ dan $E^*=E\cup E'\in\Sigma$; maka

\Centerline{$\mu
E^*\le\mu E+\mu E'<\infty$.}

\noindent Andaikan $F\in\Sigma$ dan $\mu(F\cap E^*)\le\delta^*$; maka

\Centerline{$\mu(F\cap H\cap E)\le\mu(F\cap E)\le\delta^*\le\delta$,
\quad $\mu(F\cap E')\le\delta^*$}

\noindent sehingga

\Centerline{$|(\nu+\nuprime)F|\le|\nu
F|+|\nuprime F|\le\eta+\eta\le\epsilon$,}

\Centerline{$|(c\nu)F|=|c||\nu F|\le|c|\eta\le\epsilon$,}

\Centerline{$|\nu_HF|=|\nu(F\cap H)|\le\eta\le\epsilon$.}

\noindent Karena $\epsilon$ sembarang, $\nu+\nuprime$, $c\nu$, dan
$\nu_H$ semuanya kontinu sejati.

Sekarang, dengan mengambil $H$ dari 231Eb, kita melihat bahwa
$\nu_1=\nu_H$ dan $\nu_2=-\nu_{X\setminus H}$ kontinu sejati dan tak
negatif, serta $\nu=\nu_1-\nu_2$ merupakan selisih ukuran-ukuran kontinu
sejati.
}%end of proof of 232C

\leader{232D}{Proposisi} Misalkan $(X,\Sigma,\mu)$ suatu ruang ukur,
dan $f$ suatu fungsi bernilai real yang $\mu$-terintegralkan. Untuk
$E\in\Sigma$, tetapkan $\nu E=\int_Ef$. Maka
$\nu:\Sigma\to\Bbb R$ merupakan fungsional aditif terhitung dan kontinu
sejati terhadap $\mu$, sehingga kontinu mutlak terhadap $\mu$.

\proof{ Ingat bahwa $\int_Ef=\int f\times\chi E$ terdefinisi untuk
setiap $E\in\Sigma$ (131Fa). Jadi $\nu:\Sigma\to\Bbb R$ terdefinisi
dengan baik. Jika $E$, $F\in\Sigma$ saling lepas, maka

$$\eqalign{\nu(E\cup F)
&=\int f\times\chi(E\cup F)
=\int (f\times\chi E)+(f\times\chi F)\cr
&=\int f\times\chi E+\int f\times\chi F
=\nu E+\nu F,\cr}$$

\noindent sehingga $\nu$ aditif hingga.

Sekarang 225A, tanpa menggunakan ungkapan `kontinu sejati', membuktikan
persis bahwa $\nu$ kontinu sejati terhadap $\mu$. Menurut 232Bb, $\nu$
aditif terhitung dan kontinu mutlak.
}%end of proof of 232D

\medskip

\noindent{\bf Catatan} Fungsional $E\mapsto\int_Ef$ disebut
{\bf integral tak tentu} dari $f$.


\leader{232E}{}\cmmnt{ Kini akhirnya kita siap untuk teorema tersebut.

\medskip

\noindent}{\bf Teorema Radon-Nikod\'ym}
Misalkan $(X,\Sigma,\mu)$ suatu ruang ukur dan
$\nu:\Sigma\to\Bbb R$ suatu fungsi. Maka pernyataan-pernyataan berikut
ekuivalen:

(i) terdapat fungsi $\mu$-terintegralkan $f$ sedemikian sehingga
$\nu E=\int_Ef$ untuk setiap $E\in\Sigma$;

(ii) $\nu$ aditif hingga dan kontinu sejati terhadap $\mu$.

\proof{{\bf (a)} Jika $f$ suatu fungsi bernilai real yang
$\mu$-terintegralkan dan $\nu E=\int_Ef$ untuk setiap $E\in\Sigma$,
maka 232D memberi tahu kita bahwa $\nu$ aditif hingga dan kontinu sejati.

\medskip

{\bf (b)} Untuk arah sebaliknya, andaikan $\nu$ aditif hingga dan
kontinu sejati; perhatikan bahwa (menurut 232B(a-b)) $\nu E=0$ setiap
kali $\mu E=0$. Pertama-tama, andaikan $\nu$ nonnegatif dan tak nol.

Dalam hal ini, terdapat fungsi sederhana nonnegatif $f$ sedemikian
sehingga $\int f>0$ dan $\int_Ef\le\nu E$ untuk setiap $E\in\Sigma$.
\Prf\ Misalkan $H\in\Sigma$ sedemikian sehingga $\nu H>0$; tetapkan
$\epsilon={1\over 3}\nu H>0$. Ambil $E\in\Sigma$ dan $\delta>0$
sedemikian sehingga $\mu E<\infty$ dan $\nu F\le\epsilon$ setiap kali
$F\in\Sigma$ dan $\mu(F\cap E)\le\delta$; maka
$\nu(H\setminus E)\le\epsilon$, sehingga
$\nu E\ge\nu(H\cap E)\ge2\epsilon$ dan
$\mu E\ge\mu(H\cap E)>0$. Tetapkan $\mu_EF=\mu(F\cap E)$ untuk setiap
$F\in\Sigma$; maka $\mu_E$ merupakan fungsional aditif terhitung pada
$\Sigma$. Tetapkan $\nuprime=\nu-\alpha\mu_E$, dengan
$\alpha=\epsilon/\mu E$; maka $\nuprime$ merupakan fungsional aditif
terhitung dan $\nuprime E>0$. Seperti biasa, menurut 231Eb terdapat
himpunan $G\in\Sigma$ sedemikian sehingga $\nuprime F\ge 0$ jika
$F\in\Sigma$, $F\subseteq G$, tetapi $\nuprime F\le 0$ jika
$F\in\Sigma$ dan $F\cap G=\emptyset$.
Karena $\nuprime(E\setminus G)\le 0$,

\Centerline{$0<\nuprime E\le\nuprime(E\cap G)\le\nu(E\cap G)$}

\noindent dan $\mu(E\cap G)>0$. Tetapkan
$f={\alpha}\chi(E\cap G)$; maka $f$ merupakan fungsi sederhana tak
negatif dan $\int f={\alpha}\mu(E\cap G)>0$.

Jika $F\in\Sigma$, maka $\nuprime(F\cap G)\ge 0$, yakni

\Centerline{$\nu(F\cap G)\ge\alpha\mu_E(F\cap G)
=\alpha\mu(F\cap E\cap G)=\int_Ff$.}

\noindent Jadi

\Centerline{$\nu F\ge\nu(F\cap G)\ge\int_Ff$,}

\noindent sebagaimana diperlukan.\ \Qed

\medskip

{\bf (c)} Dengan tetap mengandaikan bahwa $\nu$ merupakan fungsional
aditif nonnegatif yang kontinu sejati, misalkan $\Phi$ himpunan
fungsi-fungsi sederhana nonnegatif $f:X\to\Bbb R$ sedemikian sehingga
$\int_Ef\le\nu E$ untuk setiap $E\in\Sigma$; fungsi konstan $\tbf{0}$
termasuk dalam $\Phi$, sehingga $\Phi$ tidak kosong.

Jika $f$, $g\in\Phi$, maka $f\vee g\in\Phi$, dengan $(f\vee
g)(x)=\max(f(x),g(x))$ untuk $x\in X$. \Prf\ Tetapkan
$H=\{x:(f-g)(x)\ge 0\}\in\Sigma$; maka $f\vee g=(f\times\chi
H)+(g\times\chi(X\setminus H))$ merupakan fungsi sederhana nonnegatif,
dan untuk setiap $E\in\Sigma$,

\Centerline{$\int_Ef\vee g=\int_{E\cap H}f+\int_{E\setminus H}g
\le\nu(E\cap H)+\nu(E\setminus H)=\nu E$. \Qed}

Tetapkan

\Centerline{$\gamma=\sup\{\int f:f\in\Phi\}\le\nu X<\infty$.}

\noindent Pilih suatu barisan $\sequencen{f_n}$ dalam $\Phi$ sedemikian
sehingga $\lim_{n\to\infty}\int f_n=\gamma$. Untuk setiap $n$, tetapkan
$g_n=f_0\vee f_1\vee\ldots\vee f_n$; maka $g_n\in\Phi$ dan
$\int f_n\le\int g_n\le\gamma$ untuk setiap $n$, sehingga
$\lim_{n\to\infty}\int g_n=\gamma$. Menurut Teorema B.Levi,
$f=\lim_{n\to\infty}g_n$ terintegralkan dan $\int f=\gamma$. Perhatikan
bahwa jika $E\in\Sigma$, maka

\Centerline{$\int_Ef=\lim_{n\to\infty}\int_Eg_n\le\nu E$.}


\vthsp\Quer\ Andaikan, jika mungkin, terdapat $H\in\Sigma$ sedemikian
sehingga $\int_Hf\ne\nu H$. Tetapkan

\Centerline{$\nu_1F=\nu F-\int_Ff\ge 0$}

\noindent untuk setiap $F\in\Sigma$; maka menurut bagian (a) dari bukti ini
dan 232C, $\nu_1$ merupakan fungsional aditif hingga yang kontinu sejati,
dan kita mengandaikan bahwa $\nu_1\ne 0$. Menurut bagian (b) dari bukti ini,
terdapat fungsi sederhana nonnegatif $g$ sedemikian sehingga
$\int_Fg\le\nu_1F$ untuk setiap $F\in\Sigma$ dan $\int g>0$. Ambil
$n\in\Bbb N$ sedemikian sehingga $\int f_n+\int g>\gamma$. Maka
$f_n+g$ merupakan fungsi sederhana nonnegatif dan

\Centerline{$\int_F(f_n+g)=\int_Ff_n+\int_Fg\le\int_Ff+\int_Fg
=\nu F-\nu_1F+\int_Fg\le\nu F$}

\noindent untuk setiap $F\in\Sigma$, sehingga $f_n+g\in\Phi$, dan

\Centerline{$\gamma<\int f_n+\int g=\int f_n+g\le\gamma$,}

\noindent yang mustahil.\ \BanG\ Jadi $\int_Hf=\nu H$ untuk setiap
$H\in\Sigma$.

\medskip

{\bf (d)} Ini membuktikan teorema untuk $\nu$ nonnegatif. Untuk $\nu$
umum, kita hanya perlu mengamati bahwa $\nu$ dapat dinyatakan sebagai
$\nu_1-\nu_2$, dengan $\nu_1$ dan $\nu_2$ fungsional aditif terhitung
nonnegatif yang kontinu sejati, menurut 232C; dengan demikian,
terdapat fungsi terintegralkan $f_1$, $f_2$ sedemikian sehingga
$\nu_iF=\int_Ff_i$ untuk kedua nilai $i$ dan setiap $F\in\Sigma$.
Tentu saja, $f=f_1-f_2$ terintegralkan dan $\nu F=\int_Ff$ untuk setiap
$F\in\Sigma$. Ini menuntaskan bukti.
}%end of proof of 232E

\leader{232F}{Akibat} Misalkan $(X,\Sigma,\mu)$ suatu ruang ukur
$\sigma$-hingga dan $\nu:\Sigma\to\Bbb R$ suatu fungsi. Maka terdapat
fungsi $\mu$-terintegralkan $f$ sedemikian sehingga $\nu E=\int_Ef$
untuk setiap $E\in\Sigma$ jika dan hanya jika $\nu$ aditif terhitung
dan kontinu mutlak terhadap $\mu$.

\proof{ Gabungkan 232Bc dan 232E.
}%end of proof of 232F

\leader{232G}{Akibat} Misalkan $(X,\Sigma,\mu)$ suatu ruang ukur
hingga total dan $\nu:\Sigma\to\Bbb R$ suatu fungsi. Maka terdapat
fungsi $\mu$-terintegralkan $f$ pada $X$ sedemikian sehingga
$\nu E=\int_Ef$ untuk setiap $E\in\Sigma$ jika dan hanya jika $\nu$
aditif hingga dan kontinu mutlak terhadap $\mu$.

\proof{ Gabungkan 232Bd dan 232E.
}%end of proof of 232G

\leader{232H}{Catatan}\cmmnt{ {\bf (a)} Kebanyakan penulis menganggap
232F sudah memadai sebagai `Teorema Radon-Nikod\'ym'. Menurut pandangan
saya, persoalan mengidentifikasi integral tak tentu cukup penting untuk
membenarkan suatu analisis yang berlaku pada semua ruang ukur, sekalipun
analisis itu memerlukan konsep baru (konsep fungsional `kontinu sejati').

\header{232Hb}{\bf (b)} Saya patut memberikan contoh suatu fungsional
kontinu mutlak yang tidak kontinu sejati. Berikut ini sebuah contoh
sederhana. Misalkan $X$ sembarang himpunan tak terhitung. Misalkan
$\Sigma$ aljabar-$\sigma$ terhitung-koterhitung pada $X$, dan $\nu$
ukuran terhitung-koterhitung pada $X$ (211R). Misalkan
$\mu$ pembatasan ke $\Sigma$ dari ukuran pencacahan pada $X$. Jika $\mu E=0$,
maka $E=\emptyset$ dan $\nu E=0$, sehingga $\nu$ kontinu mutlak. Namun,
untuk setiap $E$ berukuran hingga kita mempunyai $\nu(X\setminus E)=1$,
sehingga $\nu$ tidak kontinu sejati. Lihat juga 232Xf(i).

\header{232Hc}{\bf *(c)} Dalam klasifikasi yang dikembangkan dalam Bab
21, ruang $(X,\Sigma,\mu)$ dari contoh ini agak tidak teratur; misalnya,
ruang ini tidak ditentukan secara lokal dan juga tidak dapat dilokalkan, sehingga
tidak dapat dilokalkan secara ketat, walaupun lengkap dan semihingga.
Dapatkah gejala ini terjadi dalam ruang ukur yang dapat dilokalkan secara
ketat? Di sini kita berhadapan dengan pertanyaan yang memikat. Andaikan,
dalam (b), saya menggunakan gagasan yang sama tetapi dengan
$\Sigma=\Cal PX$. Tidak ada kesulitan dalam mengonstruksi $\mu$; tetapi
dapatkah sekarang terdapat $\nu$ dengan sifat-sifat yang diperlukan,
yakni suatu fungsional aditif terhitung tak nol dari $\Cal PX$ ke
$\Bbb R$ yang bernilai nol pada semua himpunan hingga? Inilah `masalah
Banach-Ulam', yang telah saya bahas secara luas di tempat lain
({\smc Fremlin 93}), dan yang akan saya tinjau kembali dalam Jilid 5.
Pertanyaan ini disinggung lagi dalam 363S di Jilid 3.

\header{232Hd}{\bf (d)} Seusai Teorema Radon-Nikod\'ym, pertanyaan
berikut langsung muncul: untuk $\nu$ tertentu, seberapa besar variasi
yang mungkin terdapat dalam $f$ yang bersesuaian? Jawabannya cukup
langsung: dua fungsi terintegralkan $f$ dan $g$ menghasilkan integral tak
tentu yang sama jika dan hanya jika keduanya sama hampir di mana-mana
(131Hb).

\header{232He}{\bf (e)} Saya telah menyatakan Teorema Radon-Nikod\'ym
dengan fungsi terintegralkan sembarang, dengan maksud menafsirkan
`keterintegralan' dalam arti luas, sesuai dengan konvensi Jilid 1. Namun,
untuk suatu fungsional aditif terhitung kontinu sejati $\nu$, kita
dapat bertanya apakah dalam suatu pengertian terdapat fungsi terintegralkan
kanonik yang mewakilinya. Jawabannya tidak. Tetapi kita tentu tidak perlu
mengambil fungsi terintegralkan sembarang dari jenis yang dibahas dalam
Bab 12. Jika $f$ sembarang fungsi terintegralkan, terdapat himpunan
koterabaikan $E$ sedemikian sehingga $f\restr E$ terukur, dan kini kita
dapat menemukan himpunan terukur koterabaikan
$G\subseteq E\cap\dom f$; jika kita menetapkan $g(x)=f(x)$ untuk
$x\in G$ dan $0$ untuk $x\in X\setminus G$, maka $f\eae g$, sehingga
$g$ mempunyai integral tak tentu yang sama dengan $f$ (sebagaimana dicatat
dalam (d) tepat di atas), sedangkan $g$ terukur dan berdomain $X$. Dengan
demikian, kita dapat membuat penghalusan yang sepele tetapi kadang-kadang
berguna pada teorema tersebut: jika $(X,\Sigma,\mu)$ suatu ruang ukur dan
$\nu:\Sigma\to\Bbb R$ aditif hingga serta kontinu sejati terhadap $\mu$,
maka terdapat fungsi $\Sigma$-terukur dan $\mu$-terintegralkan
$g:X\to\Bbb R$ sedemikian sehingga $\int_Eg=\nu E$ untuk setiap
$E\in\Sigma$.

\header{232Hf}{\bf (f)}
Sekarang ada gunanya memperkenalkan suatu definisi umum.} %end of comment
Jika $(X,\Sigma,\mu)$ suatu ruang ukur dan $\nu$ suatu fungsional bernilai
$[-\infty,\infty]$ yang didefinisikan pada suatu keluarga subhimpunan $X$,
saya akan mengatakan bahwa suatu fungsi bernilai $[-\infty,\infty]$ yang
dilambangkan dengan $f$ dan didefinisikan pada suatu subhimpunan $X$ merupakan
{\bf turunan Radon-Nikod\'ym} dari $\nu$ terhadap $\mu$ jika
$\int_Efd\mu$ terdefinisi\cmmnt{ (dalam pengertian 214D)} dan sama dengan
$\nu E$ untuk setiap $E\in\dom\nu$. \cmmnt{Dengan demikian, fungsi-fungsi
terintegralkan yang disebut $f$ dalam 232E-232G %232E 232F 232G
semuanya merupakan `turunan Radon-Nikod\'ym'; kelak kita akan menjumpai
contoh-contoh yang kurang teratur.}

\cmmnt{Jika $\nu$ suatu ukuran dan $f$ nonnegatif,
$f$ dapat disebut {\bf fungsi kerapatan}.}

\cmmnt{\header{232Hg}{\bf (g)} Dalam seluruh pembahasan di atas, saya
mengandaikan bahwa $\nu$ didefinisikan pada seluruh domain $\Sigma$ dari
$\mu$. Namun, dalam beberapa penerapan terpenting, $\nu$ hanya didefinisikan
pada suatu aljabar-$\sigma$ yang lebih kecil, $\Tau$. Dalam hal ini kita
lazimnya berusaha menerapkan hasil-hasil yang sama dengan $\mu\restrp\Tau$
sebagai pengganti $\mu$.
}%end of comment

\leader{232I}{}{\bf Dekomposisi Lebesgue suatu fungsional aditif
terhitung: Proposisi} (a) Misalkan $(X,\Sigma,\mu)$ suatu ruang ukur dan
$\nu:\Sigma\to\Bbb R$ suatu fungsional aditif terhitung. Maka $\nu$
mempunyai representasi unik

\Centerline{$\nu=\nu_s+\nu_{ac}=\nu_s+\nu_{tc}+\nu_e$,}

\noindent dengan $\nu_s$ singular terhadap $\mu$, $\nu_{ac}$ kontinu
mutlak terhadap $\mu$, $\nu_{tc}$ kontinu sejati terhadap $\mu$, dan
$\nu_e$ kontinu mutlak terhadap $\mu$ serta bernilai nol pada setiap
himpunan berukuran hingga.

(b) Jika $X=\BbbR^r$, $\Sigma$ aljabar himpunan-himpunan Borel dalam
$\BbbR^r$, dan $\mu$ pembatasan ukuran Lebesgue ke $\Sigma$, maka $\nu$
dapat dinyatakan secara unik sebagai $\nu_p+\nu_{cs}+\nu_{ac}$, dengan
$\nu_{ac}$ kontinu mutlak terhadap $\mu$, $\nu_{cs}$ singular terhadap
$\mu$ dan bernilai nol pada himpunan-himpunan satu titik, serta
$\nu_pE=\sum_{x\in E}\nu_p\{x\}$ untuk setiap $E\in\Sigma$.

\proof{{\bf (a)(i)} Andaikan terlebih dahulu bahwa $\nu$ nonnegatif.
Dalam hal ini, tetapkan

\Centerline{$\nu_sE=\sup\{\nu(E\cap F):F\in\Sigma,\,\mu F=0\}$,}

\Centerline{$\nu_tE=\sup\{\nu(E\cap F):F\in\Sigma,\,\mu F<\infty\}$.}

\noindent Maka $\nu_s$ dan $\nu_t$ keduanya aditif terhitung.
\prooflet{\Prf\
Jelas bahwa $\nu_s\emptyset=\nu_t\emptyset=0$. Misalkan
$\sequencen{E_n}$ suatu barisan saling lepas dalam $\Sigma$ dengan
gabungan $E$. \grheada\ Jika $F\in\Sigma$ dan $\mu F=0$, maka

\Centerline{$\nu(E\cap F)=\sum_{n=0}^{\infty}\nu(E_n\cap F)
\le\sum_{n=0}^{\infty}\nu_s(E_n)$;}

\noindent karena $F$ sembarang,

\Centerline{$\nu_sE\le\sum_{n=0}^{\infty}\nu_sE_n$.}


\noindent \grheadb\ Jika $F\in\Sigma$ dan $\mu F<\infty$, maka

\Centerline{$\nu(E\cap F)=\sum_{n=0}^{\infty}\nu(E_n\cap F)
\le\sum_{n=0}^{\infty}\nu_t(E_n)$;}

\noindent karena $F$ sembarang,

\Centerline{$\nu_tE\le\sum_{n=0}^{\infty}\nu_tE_n$.}

\noindent \grheadc\ Jika $\epsilon>0$, maka (karena
$\sum_{n=0}^{\infty}\nu E_n=\nu E<\infty$) terdapat $n\in\Bbb N$
sedemikian sehingga $\sum_{k=n+1}^{\infty}\nu E_k\le\epsilon$.
Sekarang, untuk setiap $k\le n$, terdapat $F_k\in\Sigma$ sedemikian
sehingga $\mu F_k=0$ dan
$\nu(E_k\cap F_k)\ge\nu_sE_k-{{\epsilon}\over{n+1}}$. Dalam hal ini,
$F=\bigcup_{k\le n}F_k\in\Sigma$, $\mu F=0$, dan

\Centerline{$\nu_sE\ge\nu(E\cap F)
\ge\sum_{k=0}^n\nu(E_k\cap F_k)\ge\sum_{k=0}^n\nu_sE_k-\epsilon
\ge\sum_{k=0}^{\infty}\nu_sE_k-2\epsilon$,}

\noindent karena

\Centerline{$\sum_{k=n+1}^{\infty}\nu_sE_k
\le\sum_{k=n+1}^{\infty}\nu E_k\le\epsilon$.}

\noindent Karena $\epsilon$ sembarang,

\Centerline{$\nu_sE\ge\sum_{k=0}^{\infty}\nu_sE_k$.}

\noindent
\grheadd\ Demikian pula, untuk setiap $k\le n$, terdapat
$F'_k\in\Sigma$ sedemikian sehingga $\mu F'_k<\infty$ dan
$\nu(E_k\cap F'_k)\ge\nu_tE_k-{{\epsilon}\over{n+1}}$. Dalam hal ini,
$F'=\bigcup_{k\le n}F'_k\in\Sigma$, $\mu F'<\infty$, dan

\Centerline{$\nu_tE\ge\nu(E\cap F')\ge\sum_{k=0}^n\nu(E_k\cap F'_k)
\ge\sum_{k=0}^n\nu_tE_k-\epsilon
\ge\sum_{k=0}^{\infty}\nu_tE_k-2\epsilon$,}

\noindent karena

\Centerline{$\sum_{k=n+1}^{\infty}\nu_tE_k
\le\sum_{k=n+1}^{\infty}\nu E_k\le\epsilon$.}

\noindent Karena $\epsilon$ sembarang,

\Centerline{$\nu_tE\ge\sum_{k=0}^{\infty}\nu_tE_k$.}

\noindent \grheade\ Dengan menggabungkan hasil-hasil ini,
$\nu_sE=\sum_{n=0}^{\infty}\nu_sE_n$ dan
$\nu_tE=\sum_{n=0}^{\infty}\nu_tE_n$. Karena $\sequencen{E_n}$
sembarang, $\nu_s$ dan $\nu_t$ aditif terhitung.\ \Qed}


\medskip

\quad{\bf (ii)} Dengan tetap mengandaikan bahwa $\nu$ nonnegatif,
jika kita memilih suatu barisan $\sequencen{F_n}$ dalam $\Sigma$
sedemikian sehingga $\mu F_n=0$ untuk setiap $n$ dan
$\lim_{n\to\infty}\nu F_n=\nu_sX$, maka
$F^*=\bigcup_{n\in\Bbb N}F_n$ memenuhi $\mu F^*=0$ dan
$\nu F^*=\nu_sX$; dengan demikian $\nu_s(X\setminus F^*)=0$, dan
$\nu_s$ singular terhadap $\mu$ dalam pengertian 232Ac.

Perhatikan bahwa $\nu_sF=\nu F$ setiap kali $\mu F=0$. Jadi, jika kita
menuliskan $\nu_{ac}=\nu-\nu_s$, maka $\nu_{ac}$ merupakan fungsional
aditif terhitung dan $\nu_{ac}F=0$ setiap kali $\mu F=0$; yakni,
$\nu_{ac}$ kontinu mutlak terhadap $\mu$.

Jika kita menuliskan $\nu_{tc}=\nu_t-\nu_s$, maka $\nu_{tc}$ merupakan
fungsional aditif terhitung nonnegatif; $\nu_{tc}F=0$ setiap kali
$\mu F=0$, dan jika $\nu_{tc}E>0$, terdapat himpunan $F$ dengan
$\mu F<\infty$ dan $\nu_{tc}(E\cap F)>0$. Jadi, menurut 232Bb,
$\nu_{tc}$ kontinu sejati terhadap $\mu$. Tetapkan
$\nu_e=\nu-\nu_t=\nu_{ac}-\nu_{tc}$.

Dengan demikian, untuk setiap fungsional aditif terhitung nonnegatif
$\nu$, kita mempunyai representasi

\Centerline{$\nu=\nu_s+\nu_{ac}$,\quad $\nu_{ac}=\nu_{tc}+\nu_e$}

\noindent dengan $\nu_s$, $\nu_{ac}$, $\nu_{tc}$, dan $\nu_e$ semuanya
fungsional aditif terhitung nonnegatif, $\nu_s$ singular terhadap $\mu$,
$\nu_{ac}$ dan $\nu_e$ kontinu mutlak terhadap $\mu$, $\nu_{tc}$ kontinu
sejati terhadap $\mu$, dan $\nu_eF=0$ setiap kali $\mu F<\infty$.


\medskip

\quad{\bf (iii)} Untuk fungsional aditif terhitung umum
$\nu:\Sigma\to\Bbb R$, kita dapat menyatakan $\nu$ sebagai
$\nuprime-\nu\mskip0.5mu''$, dengan $\nuprime$ dan
$\nu\mskip0.5mu''$ fungsional aditif terhitung nonnegatif. Jika kita
mendefinisikan $\nuprime_s$, $\nu_s'',\ldots,\nu_e''$ seperti dalam
(i)-(ii), kita memperoleh fungsional-fungsional aditif terhitung

\Centerline{$\nu_s=\nu_s'-\nu_s''$,
\quad $\nu_{ac}=\nu_{ac}'-\nu_{ac}''$,
\quad $\nu_{tc}=\nu_{tc}'-\nu_{tc}''$,
\quad $\nu_e=\nu_e'-\nu_e''$}

\noindent sedemikian sehingga $\nu_s$ singular terhadap $\mu$ (jika
$F'$, $F''$ dipilih sedemikian sehingga

\Centerline{$\mu F'=\mu F''=\nu_s'(X\setminus F')=\nu_s''(X\setminus
F'')=0$,}

\noindent maka $\mu(F'\cup F'')=0$ dan $\nu_sE=0$ setiap kali
$E\subseteq X\setminus(F'\cup F'')$), $\nu_{ac}$ kontinu mutlak terhadap
$\mu$, $\nu_{tc}$ kontinu sejati terhadap $\mu$, dan $\nu_eF=0$ setiap
kali $\mu F<\infty$, sedangkan

\Centerline{$\nu=\nu_s+\nu_{ac}=\nu_s+\nu_{tc}+\nu_e$.}

\medskip

\quad{\bf (iv)} Selain itu, dekomposisi-dekomposisi ini unik.
\prooflet{\Prf\grheada\
Jika, misalnya, $\nu=\tilde\nu_s+\tilde\nu_{ac}$, dengan $\tilde\nu_s$
singular dan $\tilde\nu_{ac}$ kontinu mutlak terhadap $\mu$, misalkan
$F$, $\tilde F$ sedemikian sehingga $\mu F=\mu\tilde F=0$ dan
$\tilde\nu_sE=0$ setiap kali $E\cap \tilde F=\emptyset$, serta
$\nu_sE=0$ setiap kali $E\cap F=\emptyset$; maka kita mesti mempunyai

\Centerline{$\nu_{ac}(E\cap(F\cup\tilde
F))=\tilde\nu_{ac}(E\cap(F\cup\tilde F))=0$}

\noindent untuk setiap $E\in\Sigma$, sehingga

\Centerline{$\nu_sE=\nu(E\cap(F\cup\tilde
F))=\tilde\nu_sE$}

\noindent untuk setiap $E\in\Sigma$.
Dengan demikian, $\tilde\nu_s=\nu_s$ dan
$\tilde\nu_{ac}=\nu_{ac}$.
\medskip

\qquad\grheadb\ Demikian pula, jika
$\nu_{ac}=\tilde\nu_{tc}+\tilde\nu_e$, dengan $\tilde\nu_{tc}$ kontinu
sejati terhadap $\mu$ dan $\tilde\nu_eF=0$ setiap kali $\mu F<\infty$,
maka terdapat barisan-barisan $\sequencen{F_n}$ dan
$\sequencen{\tilde F_n}$ yang terdiri atas himpunan-himpunan berukuran
hingga sedemikian sehingga $\nu_{tc}F=0$ setiap kali
$F\cap\bigcup_{n\in\Bbb N}F_n=\emptyset$ dan $\tilde\nu_{tc}F=0$ setiap
kali $F\cap\bigcup_{n\in\Bbb N}\tilde F_n=\emptyset$. Tuliskan
$F^*=\bigcup_{n\in\Bbb N}(F_n\cup\tilde F_n)$; maka
$\tilde\nu_eE=\nu_eE=0$ setiap kali $E\subseteq F^*$ dan
$\tilde\nu_{tc}E=\nu_{tc}E=0$ setiap kali $E\cap F^*=\emptyset$, sehingga
$\nu_eE=\nu_{ac}(E\setminus F^*)=\tilde\nu_eE$ untuk setiap
$E\in\Sigma$, serta $\nu_e=\tilde\nu_e$ dan
$\nu_{tc}=\tilde\nu_{tc}$.\ \Qed}

\medskip

{\bf (b)} Dalam hal ini, $\mu$ $\sigma$-hingga (bdk.\ 211P), sehingga
setiap fungsional aditif terhitung yang kontinu mutlak bersifat kontinu
sejati (232Bc), dan kita selalu mempunyai $\nu_e=0$ serta
$\nu_{ac}=\nu_{tc}$. Di sisi lain, kita mengetahui bahwa semua himpunan
satu titik, dan karena itu semua himpunan terhitung, terukur. Dengan
demikian, kita mempunyai dekomposisi lebih lanjut
$\nu_s=\nu_p+\nu_{cs}$, dengan suatu himpunan terhitung
$K\subseteq\BbbR^r$ sedemikian sehingga $\nu_pE=0$ setiap kali
$E\in\Sigma$ dan $E\cap K=\emptyset$, sedangkan $\nu_{cs}$ singular
terhadap $\mu$ dan bernilai nol pada himpunan-himpunan terhitung.
\prooflet{\Prf\ (i) Jika $\nu\ge 0$, tetapkan

\Centerline{$\nu_pE=\sup\{\nu(E\cap K):K\subseteq\BbbR^r$
terhitung$\}$;}

\noindent persis seperti untuk $\nu_s$ yang ditangani dalam (a) di atas,
$\nu_p$ aditif terhitung dan terdapat himpunan terhitung
$K\subseteq\BbbR^r$ sedemikian sehingga $\nu_pE=\nu(E\cap K)$ untuk
setiap $E\in\Sigma$. (ii) Untuk $\nu$ umum, kita dapat menyatakan $\nu$
sebagai $\nuprime-\nu\mskip0.5mu''$, dengan $\nuprime$ dan
$\nu\mskip0.5mu''$ nonnegatif, serta menuliskan
$\nu_p=\nu_p'-\nu_p''$. (iii) $\nu_p$ dicirikan dengan mengatakan bahwa
terdapat himpunan terhitung $K$ sedemikian sehingga
$\nu_pE=\nu(E\cap K)$ untuk setiap $E\in\Sigma$ dan $\nu\{x\}=0$ untuk
setiap $x\in \BbbR^r\setminus K$. (iv) Jadi, jika kita menetapkan
$\nu_{cs}=\nu_s-\nu_p$, maka $\nu_{cs}$ singular terhadap $\mu$ dan
bernilai nol pada himpunan-himpunan terhitung.\ \Qed}

Sekarang, untuk setiap $E\in\Sigma$,

\Centerline{$\nu_pE=\nu(E\cap K)=\sum_{x\in K\cap E}\nu\{x\}
=\sum_{x\in E}\nu\{x\}$.}

}%end of proof of 232I

\medskip

\noindent{\bf Catatan} Representasi $\nu=\nu_p+\nu_{cs}+\nu_{ac}$
dalam (b) adalah {\bf dekomposisi Lebesgue} dari $\nu$.

\exercises{
\leader{232X}{Latihan dasar (a)}
Misalkan $(X,\Sigma,\mu)$ suatu ruang ukur dan
$\nu:\Sigma\to\Bbb R$ suatu fungsional aditif terhitung yang kontinu
mutlak terhadap $\mu$. Tunjukkan bahwa pernyataan-pernyataan berikut
ekuivalen: (i) $\nu$ kontinu sejati terhadap $\mu$; (ii) terdapat suatu
barisan $\sequencen{E_n}$ dalam $\Sigma$ sedemikian sehingga
$\mu E_n<\infty$ untuk setiap $n\in\Bbb N$ dan $\nu F=0$ setiap kali
$F\in\Sigma$ dan $F\cap\bigcup_{n\in\Bbb N}E_n=\emptyset$.

\sqheader 232Xb Misalkan $g:\Bbb R\to\Bbb R$ suatu fungsi terbatas tak
menurun dan $\mu_g$ ukuran Lebesgue-Stieltjes yang bersesuaian (114Xa).
Tunjukkan bahwa $\mu_g$ kontinu mutlak (secara ekuivalen, kontinu sejati)
terhadap ukuran Lebesgue jika dan hanya jika pembatasan $g$ pada setiap
interval tertutup terbatas bersifat kontinu mutlak dalam pengertian 225B.

\header{232Xc}{\bf (c)} Misalkan $X$ suatu himpunan dan $\Sigma$ suatu
aljabar-$\sigma$ dari subhimpunan-subhimpunan $X$; misalkan
$\nu:\Sigma\to\Bbb R$ suatu fungsional aditif terhitung. Misalkan
$\Cal I$ suatu {\bf ideal} dari $\Sigma$, yakni suatu subhimpunan dari
$\Sigma$ sedemikian sehingga ($\alpha$) $\emptyset\in\Cal I$,
($\beta$) $E\cup F\in\Cal I$ untuk semua $E$, $F\in\Cal I$, dan
($\gamma$) jika $E\in\Sigma$, $F\in\Cal I$, serta $E\subseteq F$, maka
$E\in\Cal I$. Tunjukkan bahwa $\nu$ mempunyai dekomposisi unik
$\nu=\nu_{\Cal I}+\nuprime_{\Cal I}$, dengan $\nu_{\Cal I}$ dan
$\nuprime_{\Cal I}$ fungsional aditif terhitung,
$\nuprime_{\Cal I}E=0$ untuk setiap $E\in\Cal I$, dan setiap kali
$E\in\Sigma$, $\nu_{\Cal I}E\ne 0$, terdapat $F\in\Cal I$ sedemikian
sehingga $\nu_{\Cal I}(E\cap F)\ne 0$.

\sqheader 232Xd Misalkan $X$ suatu himpunan tak kosong dan $\Sigma$
suatu aljabar-$\sigma$ dari subhimpunan-subhimpunan $X$. Tunjukkan bahwa
untuk setiap barisan $\sequencen{\nu_n}$ yang terdiri atas
fungsional-fungsional aditif terhitung pada $\Sigma$, terdapat suatu
ukuran peluang $\mu$ pada $X$ dengan domain
$\Sigma$ sedemikian sehingga setiap $\nu_n$ kontinu mutlak terhadap
$\mu$.
\Hint{mulailah dengan kasus $\nu_n\ge 0$.}

\header{232Xe}{\bf (e)} Misalkan $(X,\Sigma,\mu)$ suatu ruang ukur dan
$(X,\hat\Sigma,\hat\mu)$ pelengkapannya (212C). Misalkan
$\nu:\Sigma\to\Bbb R$ suatu fungsional aditif sedemikian sehingga
$\nu E=0$ setiap kali $\mu E=0$.
Tunjukkan bahwa $\nu$ mempunyai perluasan unik menjadi fungsional aditif
$\hat\nu:\hat\Sigma\to\Bbb R$ sedemikian sehingga $\hat\nu E=0$ setiap
kali $\hat\mu E=0$.

\spheader 232Xf Misalkan $\Cal F$ suatu ultrafilter pada $\Bbb N$ yang
memuat filter $\{\Bbb N\setminus I:I\subseteq\Bbb N$ berhingga$\}$
(2A1O). Definisikan $\nu:\Cal P\Bbb N\to\{0,1\}$ dengan
menetapkan $\nu E=1$ jika $E\in\Cal F$, dan $0$ jika
$E\in\Cal P\Bbb N\setminus\Cal F$.
(i) Misalkan $\mu_1$ ukuran pencacahan pada $\Cal P\Bbb N$. Tunjukkan
bahwa $\nu$ aditif dan kontinu mutlak terhadap $\mu_1$, tetapi tidak
kontinu sejati. (ii) Definisikan $\mu_2:\Cal P\Bbb N\to[0,1]$ dengan
menetapkan $\mu_2E=\sum_{n\in E}2^{-n-1}$. Tunjukkan bahwa $\nu$ bernilai
nol pada himpunan-himpunan yang $\mu_2$-terabaikan, tetapi tidak kontinu
mutlak terhadap $\mu_2$.

\header{232Xg}{\bf (g)} Tuliskan kembali bagian ini dalam konteks
fungsional aditif bernilai kompleks.

\spheader 232Xh Misalkan $(X,\Sigma,\mu)$ suatu ruang ukur, dan $\nu$
serta $\lambda$ fungsional aditif pada $\Sigma$, dengan $\nu$ positif dan
aditif terhitung, sehingga $(X,\Sigma,\nu)$ juga suatu ruang ukur.
(i) Tunjukkan bahwa jika $\nu$ kontinu mutlak terhadap $\mu$ dan
$\lambda$ kontinu mutlak terhadap $\nu$, maka $\lambda$ kontinu mutlak
terhadap $\mu$. (ii) Tunjukkan bahwa jika $\nu$ kontinu sejati terhadap
$\mu$ dan $\lambda$ kontinu mutlak terhadap $\nu$, maka $\lambda$
kontinu sejati terhadap $\mu$.
%232B

\leader{232Y}{Latihan lanjutan (a)} Misalkan $(X,\Sigma,\mu)$ suatu
ruang ukur dan $\nu:\Sigma\to\Bbb R$ suatu fungsional aditif hingga.
Jika $E$, $F$, $H\in\Sigma$ dan $\mu H<\infty$, tetapkan
$\rho_H(E,F)=\mu(H\cap(E\symmdiff F))$. (i) Tunjukkan bahwa $\rho_H$
merupakan pseudometrik pada $\Sigma$ (2A3Fa). (ii)
Misalkan $\frak T$ topologi pada $\Sigma$ yang dibangkitkan oleh
$\{\rho_H:H\in\Sigma,\,\mu H<\infty\}$ (2A3Fc). Tunjukkan bahwa $\nu$
kontinu terhadap $\frak T$ jika dan hanya jika ia kontinu sejati
dalam pengertian 232Ab.
($\frak T$ adalah topologi {\bf konvergensi dalam ukuran} pada $\Sigma$.)

\header{232Yb}{\bf (b)} Untuk fungsi tak menurun
$F:[a,b]\to\Bbb R$, dengan $a<b$, misalkan $\nu_F$ ukuran
Lebesgue-Stieltjes yang bersesuaian. Tunjukkan bahwa jika kita
mendefinisikan $(\nu_F)_{ac}$, dan seterusnya, terhadap ukuran Lebesgue
pada $[a,b]$, seperti dalam 232I, maka

\Centerline{$(\nu_F)_p=\nu_{F_p}$, \quad$(\nu_F)_{ac}=\nu_{F_{ac}}$,
\quad$(\nu_F)_{cs}=\nu_{F_{cs}}$,}

\noindent dengan $F_p$, $F_{cs}$, dan $F_{ac}$ didefinisikan seperti
dalam 226C.

\header{232Yc}{\bf (c)} Perluas gagasan dalam (b) ke fungsi umum $F$
yang bervariasi terbatas.

\header{232Yd}{\bf (d)} Perluas gagasan-gagasan dalam (b) dan (c) ke
interval terbuka, setengah terbuka, dan tak berbatas (bdk.\ 226Yb).

\header{232Ye}{\bf (e)} Misalkan $(X,\Sigma,\mu)$ suatu ruang ukur dan
$(X,\tilde\Sigma,\tilde\mu)$ versi c.l.d.\ yang bersesuaian (213E).
Misalkan
$\nu:\Sigma\to\Bbb R$ suatu fungsional aditif yang kontinu sejati
terhadap $\mu$. Tunjukkan bahwa $\nu$ mempunyai perluasan unik menjadi
fungsional $\tilde\nu:\tilde\Sigma\to\Bbb R$ yang kontinu sejati
terhadap $\tilde\mu$.

\header{232Yf}{\bf (f)} Misalkan $(X,\Sigma,\mu)$ suatu ruang ukur dan
$f$ suatu fungsi bernilai real yang $\mu$-terintegralkan. Tunjukkan bahwa
integral tak tentu dari $f$ merupakan satu-satunya fungsional aditif
terhitung $\nu:\Sigma\to\Bbb R$ sedemikian sehingga setiap kali
$E\in\Sigma$ dan $f(x)\in[a,b]$ untuk hampir setiap $x\in E$, berlaku
$a\mu E\le\nu E\le b\mu E$.

\header{232Yg}{\bf (g)} Katakan bahwa dua fungsional aditif terbatas
$\nu_1$, $\nu_2$ pada suatu aljabar himpunan $\Sigma$ bersifat
{\bf saling singular} jika untuk setiap $\epsilon>0$ terdapat
$H\in\Sigma$ sedemikian sehingga

\Centerline{$\sup\{|\nu_1F|:F\in\Sigma,\,F\subseteq H\}\le\epsilon$,}

\Centerline{$\sup\{|\nu_2F|:F\in\Sigma,\,F\cap H=\emptyset\}
\le\epsilon$.}

\quad {(i)} Tunjukkan bahwa $\nu_1$ dan $\nu_2$ saling singular jika dan
hanya jika, dalam bahasa 231Ya-231Yb, $|\nu_1|\wedge|\nu_2|=0$.

\quad{(ii)} Tunjukkan bahwa jika $\Sigma$ suatu aljabar-$\sigma$ dan
$\nu_1$, $\nu_2$ aditif terhitung, maka keduanya saling singular jika
dan hanya jika terdapat $H\in\Sigma$ sedemikian sehingga $\nu_1F=0$
setiap kali $F\in\Sigma$ dan $F\subseteq H$, sedangkan $\nu_2F=0$
setiap kali $F\in\Sigma$ dan $F\cap H=\emptyset$.

\quad{(iii)} Tunjukkan bahwa jika $\nu_s$, $\nu_{tc}$, dan $\nu_e$
didefinisikan dari $\nu$ dan $\mu$ seperti dalam 232I, maka setiap
pasangan dari ketiganya saling singular.

\spheader 232Yh Misalkan $(X,\Sigma,\mu)$ suatu ruang ukur dan $f$ suatu
fungsi bernilai real nonnegatif yang terintegralkan pada $X$; misalkan
$\nu$ integral tak tentu dari fungsi tersebut. Tunjukkan bahwa untuk setiap fungsi
$g:X\to\Bbb R$, $\int g\,d\nu=\int f\times g\,d\mu$, dalam arti bahwa
jika salah satu integral ini terdefinisi dalam $[-\infty,\infty]$, maka
yang lain pun demikian, dan keduanya sama. \Hint{mulailah dengan fungsi
sederhana $g$.}

\spheader 232Yi Misalkan $(X,\Sigma,\mu)$ suatu ruang ukur, $f$ suatu
fungsi terintegralkan, dan $\nu:\Sigma\to\Bbb R$ integral tak tentu dari
$f$. Tunjukkan bahwa $|\nu|$, sebagaimana didefinisikan dalam 231Ya,
merupakan integral tak tentu dari $|f|$.
%232D

\spheader 232Yj Misalkan $X$ suatu himpunan, $\Sigma$ suatu
aljabar-$\sigma$ dari subhimpunan-subhimpunan $X$, dan
$\nu:\Sigma\to\Bbb R$ suatu fungsional aditif terhitung. Tunjukkan bahwa
$\nu$ mempunyai turunan Radon-Nikod\'ym terhadap $|\nu|$ sebagaimana
didefinisikan dalam 231Ya, dan bahwa setiap turunan demikian mempunyai
modulus sama dengan $1\,\,|\nu|$-hampir di mana-mana.
%232D

\spheader 232Yk (H.K\"onig)
Misalkan $X$ suatu himpunan dan $\mu$, $\nu$ dua ukuran pada $X$ dengan
domain yang sama, $\Sigma$. Untuk $\alpha\ge 0$ dan $E\in\Sigma$,
tetapkan $(\alpha\mu\wedge\nu)(E)
=\inf\{\alpha\mu(E\cap F)+\nu(E\setminus F):F\in\Sigma\}$ (bdk.\
112Ya\formerly{1{}12Yb}).
Tunjukkan bahwa pernyataan-pernyataan berikut ekuivalen: (i) $\nu E=0$
setiap kali $\mu E=0$; (ii)
$\sup_{\alpha\ge 0}(\alpha\mu\wedge\nu)(E)=\nu E$ untuk setiap
$E\in\Sigma$.
}%end of exercises

\endnotes{
\Notesheader{232} Teorema Radon-Nikod\'ym harus termasuk dalam setiap
daftar enam teorema terpenting teori ukuran; bukan hanya teoremanya
sendiri, melainkan juga teknik-teknik yang diperlukan untuk membuktikannya,
merupakan inti bidang ini. Dalam buku saya {\smc Fremlin 74}, saya
membahas berbagai versi yang kurang lebih abstrak dari teorema dan metode
tersebut; sebagian di antaranya akan saya tinjau kembali dalam \S\S327 dan
365 di jilid berikutnya.

Sebagaimana saya menyajikannya di sini, inti pembuktian terbagi antara
231E dan 232E. Saya pikir kita dapat membedakan unsur-unsur berikut.
Misalkan $\nu$ suatu fungsional aditif terhitung.

\inset{(i) $\nu$ terbatas (231Ea).}

\inset{(ii) $\nu$ dapat dinyatakan sebagai selisih fungsional-fungsional
nonnegatif (231F).}

\noindent (Saya memberikannya sebagai suatu akibat dari 231Eb, tetapi
hasil ini juga dapat dibuktikan dengan metode yang lebih sederhana,
seperti dalam 231Ya.)

\inset{(iii) Jika $\nu>0$, terdapat fungsi terintegralkan $f$ sedemikian
sehingga $0<\nu_f\le\nu$,}

\noindent dengan menuliskan $\nu_f$ untuk integral tak tentu dari $f$.
(Pada titik inilah kita benar-benar memerlukan dekomposisi Hahn 231Eb.)

\inset{(iv) Himpunan $\Psi=\{f:\nu_f\le\nu\}$ tertutup terhadap
supremum terhitung, sehingga terdapat $f\in\Psi$ yang memaksimumkan
$\int f$.}

\noindent (Dalam bagian (b) dari pembuktian 232E, saya berbicara tentang
fungsi sederhana; tetapi ini semata-mata untuk menyederhanakan perincian
teknis, dan argumen yang sama berlaku jika kita menerapkannya pada $\Psi$
alih-alih $\Phi$. Perhatikan penggunaan Teorema B.Levi di sini.)

\inset{(v) Ambil $f$ dari (iv) dan gunakan (iii) untuk menunjukkan bahwa
$\nu-\nu_f=0$.}

\noindent Setiap langkah (i)-(iv) memerlukan suatu gagasan yang tidak
sepele, dan pentingnya teorema ini tidak hanya terletak pada konsekuensi
langsungnya yang mengagumkan dalam sisa bab ini dan di tempat lain,
melainkan juga pada keserbagunaan dan kekuatan gagasan-gagasan tersebut.

Saya memperkenalkan gagasan fungsional `kontinu sejati' untuk memberikan
uraian yang cukup langsung mengenai kedudukan Teorema Radon-Nikod\'ym
dalam ruang ukur yang tidak $\sigma$-hingga. Tentu saja, pokok utamanya
ialah bahwa suatu fungsional kontinu sejati, seperti integral tak tentu,
harus terpusat pada bagian ruang yang $\sigma$-hingga (232Xa), sehingga
232E, sebagaimana dinyatakan, dapat dengan mudah diturunkan dari bentuk
baku 232F. Saya memberanikan diri menggunakan kata `sejati' dalam konteks
ini karena jenis kekontinuan ini memang bersesuaian dengan suatu konsep
topologis (232Ya).

Ada kemungkinan jebakan dalam definisi fungsional `kontinu mutlak' yang
saya berikan. Banyak penulis menggunakan syarat 232Ba sebagai definisi,
dengan mengatakan bahwa $\nu$ kontinu mutlak terhadap $\mu$ jika
$\nu E=0$ setiap kali $\mu E=0$. Untuk fungsional aditif terhitung,
syarat ini sama dengan perumusan $\epsilon$-$\delta$ dalam 232Aa; tetapi
untuk fungsional aditif lain, hal itu tidak harus demikian (232Xf(ii)).
Pada umumnya perbedaan tersebut tidak penting, tetapi saya mencatat bahwa
dalam 232Bd perbedaan itu sangat penting, karena di sana $\nu$ tidak
diandaikan aditif terhitung.

Dalam 232I saya menguraikan salah satu dari banyak cara untuk menguraikan
suatu fungsional aditif terhitung menjadi bagian-bagian saling singular
dengan sifat-sifat khusus. Dalam 231Yf-231Yg saya telah menyarankan suatu
metode untuk menguraikan fungsional aditif menjadi jumlah suatu bagian
aditif terhitung dan suatu bagian `aditif hingga murni'. Semua hasil ini
mempunyai ungkapan alami dalam kerangka ruang linear terurut dari
fungsional-fungsional aditif terbatas pada suatu aljabar (231Yc).
}%end of notes

\discrpage
